\section{More Details of Experiment} \label{section_c: Experiments}
\subsection{Significance of New GPT-4 Based Metric} \label{appendix_c: Evaluation Metrics}
\noindent\textbf{GPT-4 prompt design}\quad 
For each of the three tasks, we designed three prompts for scoring. \textit{For Detailed Description Task (H2, C3, and M2)}, we designed the prompt in five areas: text length, text content variation, text detail variation, logical text description, and linguistic ability.
For \textit{Counter-intuitiveness Reasoning Task (H3, C3, and M3)}, we designed the prompt in six areas: expressiveness of language, the logic of response, the common sense of response, understanding of counter-intuition, differences in text detail, and length of text.
For \textit{Title Task (H4 and C4)}, we used the description, comprehension, and title of the manually annotated video as a reference to score the new title. Each Prompt can be seen in Figure \ref{fig:prompt-eval-title}, \ref{fig:prompt-eval-des}, and \ref{fig:prompt-eval-exp}.

\begin{figure}[htbp]
    \centering
    \includegraphics[width=\linewidth]{figures_appendix/prompt-eval-title.pdf}
    \vspace{-0.6cm}
    \caption{\textbf{GPT-4 prompt for Title Evaluation.} We give GPT prompt as system, and then input a description, explanation, and two titles, the first one is our annotation, and the second one is the output of the model, and GPT-4 will evaluate the similarity between these two texts according to the prompt's requirements.}
    \vspace{-0.4cm}
    \label{fig:prompt-eval-title}
\end{figure}
\begin{figure}[htbp]
    \centering
    \includegraphics[width=\linewidth]{figures_appendix/prompt-eval-des.pdf}
    \vspace{-0.6cm}
    \caption{\textbf{GPT-4 prompt for Description Evaluation.} We give GPT-4 prompt as system, and then input two descriptions, the first one is our annotation, and the second one is the output of the model, and GPT-4 will evaluate the similarity between these two texts according to the prompt's requirements.}
    \vspace{-0.4cm}
    \label{fig:prompt-eval-des}
\end{figure}
\begin{figure}[htbp]
    \centering
    \includegraphics[width=\linewidth]{figures_appendix/prompt-eval-exp.pdf}
    \vspace{-0.6cm}
    \caption{\textbf{GPT-4 prompt for Explanation Evaluation.} We give GPT prompt as system, and then input two explanations, the first one is our annotation, and the second one is the output of the model, and GPT-4 will evaluate the similarity between these two texts according to the prompt's requirements.}
    \vspace{-0.4cm}
    \label{fig:prompt-eval-exp}
\end{figure}

\noindent\textbf{Comparison between GPT-4 and Traditional Metrics}\quad
The principle of the traditional metrics is relatively simple. The traditional metric principles are as follows.\\
\noindent\textbf{BLEU-4}\quad BLEU, full name is Bilingual Evaluation Understudy, is a commonly used machine translation evaluation metric. It evaluates how good a machine translation result is by comparing how well it matches the N-gram of one or more human-translated reference results, which is a sequence of N consecutive words. BLEU-4, i.e., evaluates how well a combination of two words (i.e., a binary) matches. BLEU introduces a correction factor, Brevity Penalty (BP), to avoid this problem, which penalizes machine translation if the result is shorter than the reference translation. The BLEU score is the geometric mean of the individual N-gram accuracy multiplied by the shortness penalty. That is, the score of BLEU takes into account the precision and length of the translation result.\\
\noindent\textbf{ROUGE-L}\quad
ROUGE, known as Recall-Oriented Understudy for Gisting Evaluation, is a commonly used evaluation method for tasks such as automatic digesting and machine translation. ROUGE is mainly evaluated by comparing the overlap between the generated abstracts and the reference abstracts. Among them, ROUGE-L is an important variant of ROUGE, where L stands for Longest Common Subsequence (LCS), i.e., the longest common subsequence. Unlike the n-gram, the longest common subsequence does not require consecutive occurrences of items in the sequence.
\\
\noindent\textbf{CIDEr}\quad
CIDEr, known as Consensus-based Image Description Evaluation, is an evaluation metric specifically designed for evaluating image description (Image Captioning) tasks. The main advantage of CIDEr is that it can capture more detailed information because it uses TF-IDF weights to emphasize n-grams that occur frequently in manual annotation but are not common in all image descriptions.\\
\noindent\textbf{BLEURT}\quad
BLEURT, full name Bilingual Evaluation Understudy with Representations from Transformers, is an evaluation method based on the transformer model, specifically for evaluating the output quality of machine translation and natural language generation tasks. Unlike traditional evaluation metrics such as BLEU and ROUGE, BLEURT does not directly compare the n-gram match between generated text and reference text but uses pre-trained language models (e.g., BERT) to understand the semantic information of text. the advantage of BLEURT is that it can capture the deep semantic information of text, and can address some problems that are difficult to be handled by traditional evaluation metrics (e.g. synonym substitution, utterance rearrangement) giving reasonable evaluation.\\
It can be seen that the principles of the above traditional NLG metric are relatively simple and lack the ability to understand and evaluate the text with complex logic. In contrast, GPT-4, as an LLM, has a high ability to understand the text and can understand the text first and evaluate the similarity through prompt design.

\noindent\textbf{Instability of GPT-4 responses}\quad
During the experimental process using GPT-4 as the metric, we found that in a few cases, GPT-4 would provide different evaluations for the same content. In such cases, we would perform multiple evaluations and take the average. Figure \ref{fig:gpt-dif} shows an example of GPT-4 generating different responses to the same content.

\begin{figure}[htbp]
    \centering
    \includegraphics[width=\linewidth]{figures_appendix/GPT-4-dif.pdf}
    \vspace{-0.6cm}
    \caption{\textbf{GPT-4 generates different responses to the same content.} User Input is the two sentences input to GPT-4, which is used to judge the similarity. The GPT-4 Output 1 and GPT-4 Output 2 below are the answers given by GPT-4 in two different conversations, and you can see that the two assessments are different.}
    \vspace{-0.4cm}
    \label{fig:gpt-dif}
\end{figure}

\subsection{On Failure of the Timestamp Localization Task} \label{appendix_c: Failure of t1}
There is no model that has the ability to answer the Timestamp Localization Task. Figure \ref{fig:T1} shows exactly what is 
answered for each model.

\begin{figure}[htbp]
    \centering
    \includegraphics[width=\linewidth]{figures_appendix/Task_1.pdf}
    \vspace{-0.6cm}
    \caption{\textbf{Model responses on Timestamp Localization Task.} As can be seen, the model gives answers to the Timestamp Localization Task still focusing on describing the video content and still does not answer the specific time period when asked to type the number of frames and seconds.}
    \vspace{-0.4cm}
    \label{fig:T1}
\end{figure}

\subsection{Implementation Details} \label{appendix_c: Exp_result}
\noindent\textbf{mPLUG}\quad
mPLUG is a multi-modal system employing independent image and text encoders, a cross-modal network, and a text generation decoder, which is trained through prefix language modeling loss to generate captions from connected image and prefix sub-sequence representations.\\
\noindent\textbf{GIT}\quad 
GIT is a system with an image encoder and a text decoder; it processes multiple video frames independently, adds learnable temporal embeddings before concatenation, uses a contrastively pre-trained model for image encoding, and employs a transformer module for text prediction. We used the 14M version and used two models, GIT\_LARGE\_VATEX and GIT\_LARGE\_MSRVTT, which were fine-tuned on the video captioning task for the VATEX and MSRVTT datasets, respectively.\\
\noindent\textbf{VideoChat}\quad
VideoChat, specifically the VideoChat-13B version, is an end-to-end system for video comprehension that combines pre-trained models. It utilizes QFormer to generate video embeddings and then employs LLAMA-13B for multimodal understanding and outputs video text descriptions with timestamps. In the experiment, we used VideoChat-13B with the hyperparameters: beam search number = 1, temperature = 1, video segments = 8, and token = 512.\\
\noindent\textbf{Video-ChatGPT}\quad
Video-ChatGPT is a vision-language model with a video encoder and LLM. It generates answers using video embeddings and benefits from a data-centric, human-assisted annotation framework for high-quality video instructional data.
In the experiment, we used Video-ChatGPT-7B with its hyperparameter: temperature = 0.2, and token = 512.\\
\noindent\textbf{Otter}\quad
The Otter model employs the OpenFlamingo training paradigm, utilizing pre-trained encoders for language (LLaMA-7B) and vision (CLIP ViT-L/14). In the fine-tuning process, Otter prioritizes the Perceiver resampler module while keeping the encoders frozen. In the training stage, we finetuned Otter on Dense Caption and FunQA for a total of 3 epochs each. In the experiment, we used two versions of Otter with the same hyperparameters: beam search number = 3,  size of no-repeat-ngram = 0.2, and token = 256.

All of our experiments took five to six days on a single NVIDIA A40 48GB GPU.

\subsection{More Examples}
Figure \ref{fig:more-model-res} shows the responses of different models on CreativeQA and MagicQA.

\begin{figure}[htbp]
    \centering
    \includegraphics[width=\linewidth]{figures_appendix/More_Examples.pdf}
    \vspace{-0.6cm}
    \caption{\textbf{Model responses on CreativeQA and MagicQA. } For the description of the Creative video example, only VideoChat gives the key point of the bicycle, but its description also has many errors and omissions, and the remaining two models do not identify the bicycle. In the explanation task, the responses of all three models fail to clearly explain the creativity of this imitation performance. For the Magic video example, all three models perform very poorly in description and explanation, basically only answering the phone and the straw, but lacking the description and explanation of the magic effect.}
    \vspace{-0.4cm}
    \label{fig:more-model-res}
\end{figure}

